\cvsection{Expériences professionnelles}

\begin{cventries}
  \cventry
    {Stagiaire ingénieur de recherche en IA}
    {INRIA}
    {Bordeaux, France}
    {Mars 2026 -- Septembre 2026}
    {
      \begin{cvitems}
        \item {Conception et développement en autonomie, avec le suivi de mes encadrants, d’une application répondant au besoin d’aider les chercheurs à trouver des articles scientifiques.}
        \item {Construction d’un agent multi-étapes sans framework d’orchestration afin de maîtriser ses outils, son état, sa mémoire et ses enchaînements de tâches.}
        \item {Intégration de l’API OpenAI, conception de prompts et sélection des messages par l’utilisateur pour contrôler le contexte transmis au LLM.}
        \item {Intégration des composants applicatifs avec un backend Python/FastAPI, des API REST documentées avec OpenAPI, MongoDB et une interface React/Next.js.}
        \item {Mise en œuvre d’un pipeline RAG avec découpage des documents, embeddings, indexation dans ChromaDB et recherche sémantique.}
        \item {Gestion d’une application web conteneurisée utilisant des LLM installés localement sur un Mac Studio, avec installation et configuration des modèles.}
        \item {Utilisation de vLLM pour lancer et servir des modèles sur GPU NVIDIA.}
        \item {Comparaison de Qwen, Gemma et GLM et de plusieurs quantifications sur infrastructure HPC avec Slurm selon la qualité, la latence, le TTFT et la taille du contexte.}
        \item {Définition du protocole expérimental, analyse des résultats et rédaction d’un article scientifique présentant les résultats et les limites des modèles.}
        \item {Évaluation d’OpenClaw dans un environnement Docker isolé et analyse des risques liés à l’accès autonome aux outils et aux fichiers.}
        \item {Conteneurisation de composants avec Docker et mise en place d’un pipeline CI/CD sur GitHub.}
      \end{cvitems}
    }

  \cventry
    {Stagiaire de recherche en apprentissage par renforcement}
    {Université Karunya}
    {Coimbatore, Inde}
    {Juin 2025 -- Septembre 2025}
    {
      \begin{cvitems}
        \item {Développement d’un modèle d’apprentissage par renforcement et de son environnement d’entraînement pour la navigation autonome d’un robot hospitalier.}
        \item {Déploiement et optimisation du logiciel et du modèle sur un Raspberry Pi afin de le tester en conditions réelles.}
      \end{cvitems}
    }
\end{cventries}
